\scanpage{43}
$Z$ est une structure quasi-régulière sur $S$ (identifiée p.ex. à un ss-groupe cyclique de $\unsure{\mathfrak S_S|\mathcal S_S|G_S}$ d'ordre $n=\operatorname{card}(S)$), tels que pour une (donc tte) structure polygonale $\omega$ sur $S$ compatible avec $Z$ (i.e. telle que $\omega\subset Z$), $(E,S,\omega)$ soit un polygone régulier. Il revient au $=$ de dire que tout $u\in\mathrm{Aff}^{\pm}(S,Z)$ se prolonge de façon unique en un autom. affine de $E$. [NB On ne dit rien sur les autres automorphismes de $\mathrm{Aff}(S,Z)$ \deleted{\unsure{ca}}
% slanted
\margin{Question à élucider}
(il y en a bien d'autres, sauf dans le cas où $\unsure{\varphi(n)|\phi(n)}=2$ \inserted{nb d'Euler} i.e. $n=\deleted{2,}3,4$ ou $6$)] On peut encore dire que \inserted{la donnée d'}un \deleted{\unsure{gr|un}} syst. rég. de $n$ sommets dans $E$ équivaut à la donnée d'une orbite de $\Psi_n\inserted{=(\mathbb Z/n\mathbb Z)^{\unsure{*|\times}}/\{\pm1\}}$ opérant sur l'ens de tous les $n$-polygônes réguliers dans $E$.
% archive sheet number: 42
